% Unicode corpus (M3): polyglossia with Greek and Russian, Libertinus Serif
% by file name.
\documentclass{article}
\usepackage{fontspec}
\setmainfont{LibertinusSerif}[Extension=.otf, UprightFont=*-Regular,
  BoldFont=*-Bold, ItalicFont=*-Italic, BoldItalicFont=*-BoldItalic]
\usepackage{polyglossia}
\setdefaultlanguage{english}
\setotherlanguages{greek,russian}
\begin{document}
\section{Greek}
\begin{greek}
Η γλώσσα είναι ελληνική και έχει τονισμένα φωνήεντα. \textbf{Έντονα} και
\textit{πλάγια} γράμματα. Σήμερα είναι \today.
\end{greek}
\section{Russian}
\begin{russian}
Съешь же ещё этих мягких французских булок, да выпей чаю. Электрификация,
достопримечательность, \textbf{полужирный} и \textit{курсив}, \today.
\end{russian}
\end{document}
